% audit_patches.tex — replacement snippets for p3_article.tex v6, produced by the referee audit (see audit_report.md).
% Each patch: a comment block "ORIGINAL (p3_article.tex v6, line L):" quoting the exact text to replace, then the replacement.
% Patches are independent; apply with a text editor. Nothing in the manuscript has been modified by the audit.
% Uses the manuscript's macros (\Pthree, \Dzero, \Sym, \ctype, \Cent, \supp, \ZZ, \TR).

% =====================================================================================================================
% P-B1  Corollary cor:lambda, third clause (line 389)
% ORIGINAL: "it is polynomial when all parts of $\lambda$ are at most $2$ (then $x^{2}=1$, so $x=ca^{-1}$ and one checks
%   $\ctype(ca^{-1})=\lambda$), and in time $N^{O(s)}$ when $\lambda$ has at most $s$ parts larger than $1$."
% =====================================================================================================================
it is polynomial when all parts of $\lambda$ are at most $2$ (then $x^{2}=1$, so $x=ca^{-1}$ and one checks
$\ctype(ca^{-1})=\lambda$), and in time $N^{O(s)}$ when the parts of $\lambda$ larger than $1$ sum to at most $s$ (then
$|\supp x|\le s$, and the at most $N^{s}$ candidates are checked directly); by Lemma~\ref{lem:support} the bound $N^{O(s)}$ also
holds for every $\lambda$ when $|\supp a|\le s$.

% =====================================================================================================================
% P-B2  Merge prop:moment into thm:firstmoment.  Delete lines 142–151 (the whole \begin{prop}[first-moment identity] … \end{proof}),
% and replace the statement of thm:firstmoment (lines 161–168) and the sentence after it (line 176) as follows.
% ORIGINAL (line 161–163): "\begin{thm}[first moment, closed form]\label{thm:firstmoment}
%   For $a$ of cycle type $\alpha$ and any class $\gamma$ of $\Sym(N)$,
%   \[ \sum_{c\in\gamma}\#\mathrm{Sol}(a,c)=\#\{x:\ctype(ax^{3})=\gamma\}=\sum_{\nu}\rho_3(\nu)\,K(\gamma,\nu;\alpha), \]"
% =====================================================================================================================
\begin{thm}[first moment]\label{thm:firstmoment}\label{prop:moment}
Fix $a\in\Sym(N)$ of cycle type $\alpha$ and a conjugacy class $\gamma$. Then
\[ \sum_{c\in\gamma}\#\{x\in\Sym(N):\ x\,a\,x^{2}=c\}\;=\;\#\{x\in\Sym(N):\ \ctype(a\,x^{3})=\gamma\}\;=\;\sum_{\nu}\rho_3(\nu)\,K(\gamma,\nu;\alpha), \]
the sum over the conjugacy classes $\nu$ of $\Sym(N)$, where
% (keep lines 164–168 unchanged: definition of K and rho_3)
% ORIGINAL (line 176): "The identity was checked against the exhaustive tables on all $1813$ pairs of classes $(\alpha,\gamma)$ for $N\le9$.
%   It gives the average of the solution count over a block of pairs in closed form; inside a block the counts are approximately Poisson
%   with that mean, and nothing in the identity constrains an individual pair."
The right-hand side depends only on $\alpha$ and $\gamma$. Its evaluation is polynomial in the number $p(N)=e^{\Theta(\sqrt N)}$ of
classes once the character table of $\Sym(N)$ is available, but we know no algorithm polynomial in $N$; the identity is a statement about
averages, not an algorithm. It was checked against the exhaustive tables on all $1813$ pairs of classes $(\alpha,\gamma)$ for $3\le N\le9$
(the first equality alone on the $121$, $225$ and $484$ blocks at $N=6,7,8$). It gives the average of the solution count over a block of
pairs in closed form; inside a block the counts are approximately Poisson with that mean, and nothing in the identity constrains an
individual pair.
% Also replace, in the intro results list (line 52), "The first moment has a closed form in the class algebra (Theorem~\ref{thm:firstmoment})."
% by: "and has a closed form in the class algebra (Theorem~\ref{thm:firstmoment})."  — and in the abstract keep "an exact first-moment identity".

% =====================================================================================================================
% P-B3  Lemma lem:support: full proof.
% ORIGINAL (lines 95–97): "\begin{proof}[Proof sketch]
%   By Lemma~\ref{lem:conj}, $c=bx^{3}$ with $b=xax^{-1}$, and $b$ is determined by the images $x(p)$, $p\in\supp(a)$, of which there are
%   at most $N^{s}$ choices. … All of this is checked in polynomial time. (Verified exhaustively for $s=2$, $N=6$: $30/30$ instances.)
%   \end{proof}"
% =====================================================================================================================
\begin{proof}
Let $P=\supp(a)$, $|P|=s$. A solution $x$ restricts to an injection $f=x|_P\colon P\to[N]$; there are at most $N^{s}$ injections and we
enumerate them. Fix $f$. Since $a$ permutes $P$, the permutation $b_f$ with $b_f(f(p))=f(a(p))$ for $p\in P$ and $b_f(j)=j$ for
$j\notin f(P)$ is well defined, and $xax^{-1}=b_f$ for every $x$ extending $f$. By Lemma~\ref{lem:conj}, such an $x$ solves $xax^{2}=c$ iff
$x^{3}=t_f:=b_f^{-1}c$. It remains to decide in time $N^{O(s)}$ whether a permutation $t$ has a cube root extending an injection $f$
defined on $s$ points.

Let $\mathcal C$ be the set of cycles of $t$. A cube root $x$ of $t$ commutes with $t$, hence induces a length-preserving permutation
$\pi$ of $\mathcal C$, and $x$ is determined on a cycle $C$ by the image of one point, $x(t^{i}z)=t^{i}x(z)$. Since $x^{3}=t$, every cycle
of $\pi$ has length $1$ or $3$: a fixed cycle $C$ of $\pi$ has length prime to $3$ and $x|_C=t^{k}|_C$ with $3k\equiv1\pmod{|C|}$; on a
$3$-cycle $(C_1\,C_2\,C_3)$ of $\pi$ the three cycles have a common length $m$ and $x$ is determined by $x(z_1)\in C_2$ and $x(x(z_1))\in C_3$
for one $z_1\in C_1$, after which $x(x(x(z_1)))=t(z_1)$ ($2m^{2}$ possibilities for the triple). Conversely every length-preserving $\pi$ with
cycles of length $1$ (only on cycles of $t$ of length prime to $3$) or $3$, together with such local data, defines a cube root of $t$.

Let $\mathcal C_P$ be the set of cycles of $t$ meeting $P$; $|\mathcal C_P|\le s$. The constraint $x|_P=f$ determines $x$, hence $\pi$,
on every $C\in\mathcal C_P$ (if two constraints on one cycle disagree, or $|\pi(C)|\ne|C|$, the injection $f$ is rejected). We complete the
$\pi$-cycle through each $C\in\mathcal C_P$: if $\pi(C)=C$ we check $x|_C=t^{k}|_C$; if $\pi(C)=C'\ne C$ and $\pi(C')$ is already determined,
$\pi(\pi(C'))=C$ and $x$ on $\pi(C')$ are forced by $x^{3}=t$ and are checked; otherwise we choose $\pi(C')=C''$ among the at most $N$ cycles
of length $|C|$ not yet used and the image $x(z')\in C''$ of one point $z'\in C'$ (at most $N$ choices), after which $x$ on $C''$ is forced.
Each cycle of $\mathcal C_P$ thus costs at most $N^{2}$ choices, and at most $N^{2s}$ partial maps $x_0$ have to be examined; each is defined
on a $t$-invariant set $\Omega_0\supseteq P$ and satisfies $x_0^{3}=t|_{\Omega_0}$ and $x_0|_P=f$. Such an $x_0$ extends to a cube root of $t$
iff the cycles of $t$ outside $\Omega_0$ can be grouped into singletons of length prime to $3$ and triples of equal length, i.e.\ iff for
every length $m$ divisible by $3$ the number of those cycles of length $m$ is divisible by $3$ (cycles of length prime to $3$ may always be
singletons); when this holds, any grouping with any phases completes $x_0$. The procedure takes $N^{s}\cdot N^{2s}\cdot N^{O(1)}=N^{O(s)}$
steps and produces a solution whenever one exists.
\end{proof}
The algorithm of this proof was checked against brute force on $360$ instances with $N\in\{5,6,7\}$ and $|\supp a|\in\{0,2,3,4\}$
(\texttt{audit\_checks.py}), in addition to the earlier $30/30$ instances with $s=2$, $N=6$.

% =====================================================================================================================
% P-M1  Theorem thm:transposition: rewritten proof.
% ORIGINAL (lines 127–138): "\begin{proof}
%   By Proposition~\ref{prop:dual} count the $y$ with $y\,s\,y^{2}=a$, $s=(0\ d)$. … and $ysy^{2}=(ysy^{-1})y^{3}=(u\ v)(u\ v)a=a$.
%   \end{proof}"
% =====================================================================================================================
\begin{proof}
By Proposition~\ref{prop:dual} it suffices to count the $y$ with $y\,s\,y^{2}=a$, $s=(0\ d)$. Fix such a $y$ and put $u=y(0)$, $v=y(d)$,
so that $ysy^{-1}=(u\ v)$ and, by Lemma~\ref{lem:conj}, $y^{3}=(u\ v)^{-1}a=(u\ v)\,a$. The permutation $(u\ v)\,a$ maps $j\mapsto j+1$
except $u-1\mapsto v$ and $v-1\mapsto u$; it has exactly two cycles, the arcs $C_1=(u,u+1,\dots,v-1)$ of length $\ell=(v-u)\bmod N\in[1,N-1]$
and $C_2=(v,v+1,\dots,u-1)$ of length $N-\ell$, on each of which it acts as $j\mapsto j+1$.

Since $y$ commutes with $y^{3}$, it permutes the cycles of $y^{3}$. A cycle of $y$ of length divisible by $3$ would cube to three cycles,
so every cycle of $y$ has length prime to $3$ and cubes to one cycle of the same length: the cycles of $y$ are $C_1$ and $C_2$ as sets
(also when $\ell=N-\ell$), $3\nmid\ell$, $3\nmid N-\ell$, and on $C_i$ the permutation $y$ is the unique cube root of the rotation $+1$ in
the cyclic group generated by that rotation, i.e.\ $y$ acts on $C_i$ as $+e_i$ along the arc, $3e_i\equiv1\pmod{|C_i|}$ ($y$ fixes the
point if $|C_i|=1$). Consequently $0\in C_1$, because $y(0)=u\in C_1$, and $d\in C_2$, because $y(d)=v\in C_2$; in particular $\ell$ is the
length of the cycle of $y$ through $0$.

Write $0=u+\pi_1$ and $d=v+\pi_2$ with $\pi_1\in[0,\ell-1]$, $\pi_2\in[0,N-\ell-1]$ the positions of $0$ on $C_1$ and of $d$ on $C_2$.
Then $y(0)=u$ reads $\pi_1+e_1\equiv0\pmod\ell$, i.e.\ $\pi_1=\varepsilon(\ell)$, and $y(d)=v$ reads $\pi_2=\varepsilon(N-\ell)$. Hence
$u\equiv-\varepsilon(\ell)$, $v=u+\ell$ and $d\equiv v+\varepsilon(N-\ell)\equiv\ell+\varepsilon(N-\ell)-\varepsilon(\ell)\pmod N$, so $\ell$
is admissible, and $\ell$ determines $u$, $v$ and $y$: the map $y\mapsto\ell$ is injective. Conversely, for an admissible $\ell$ put
$u=-\varepsilon(\ell)\bmod N$, $v=u+\ell$, and let $y$ act as $+e_i$ on the two arcs of $(u\ v)\,a$. Then $y^{3}=(u\ v)\,a$; $0$ sits at
position $\varepsilon(\ell)$ on $C_1$ and $d$ at position $\varepsilon(N-\ell)$ on $C_2$ by admissibility, so $y(0)=u$, $y(d)=v$,
$ysy^{-1}=(u\ v)$ and $ysy^{2}=(ysy^{-1})\,y^{3}=(u\ v)(u\ v)\,a=a$. The solutions $y$ are therefore in bijection with the admissible $\ell$.
\end{proof}

% =====================================================================================================================
% P-M2  Lemma lem:affinep, proof (lines 286–293).
% ORIGINAL: "whose slope is the product of the $p$ slopes and whose translation part is a linear form in $t_{i_1},\dots,t_{i_p}$; the slope
%   condition is necessary, and when it holds the translation part vanishes for a $(p-1)$-dimensional affine family of choices of $t$ (one
%   linear condition on $p$ free parameters), so a solution exists on that cycle. On a fixed point $i$, $x|_{B_i}$ is $z\mapsto m_iz+t_i$ and
%   $x^{p}=1$ there iff $m_i^{p}\equiv1$ (and, when $m_i\ne1$, $t_i$ is then arbitrary; when $m_i=1$, $t_i=0$)."
% =====================================================================================================================
whose slope is the product of the $p$ slopes and whose translation part is a linear form in $t_{i_1},\dots,t_{i_p}$ in which the last
translation applied, $t_{i_p}$, has coefficient $1$; the slope condition is necessary, and when it holds the translation part vanishes for
exactly $L^{p-1}$ choices of $(t_{i_1},\dots,t_{i_p})$ (solve the one non-degenerate linear condition for $t_{i_p}$), so a solution exists
on that cycle. On a fixed point $i$, $x|_{B_i}$ is $z\mapsto m_iz+t_i$ and $x^{p}(z)=m_i^{p}z+(1+m_i+\dots+m_i^{p-1})t_i$; this is the
identity for some $t_i$ iff $m_i^{p}\equiv1\pmod L$, since $t_i=0$ then works. (For prime $L$ and $m_i\neq1$ every $t_i$ works, because
$1+m_i+\dots+m_i^{p-1}=(m_i^{p}-1)/(m_i-1)=0$; for composite $L$ this may fail, e.g.\ $L=9$, $p=3$, $m_i=4$ allows only $t_i\in\{0,3,6\}$.)

% =====================================================================================================================
% P-M3a  Proof of Theorem thm:main, first sentences (line 249).
% ORIGINAL: "Given an N3DM instance, put $M=3B$ and let $L$ be the least prime with $L\equiv2\pmod3$ and $T:=L-1>12B$; by the prime number
%   theorem in arithmetic progressions $L=O(B)$, and for $B\le10^{4}$ one checks directly that $L\le2(12B+2)$; it is found by trial division
%   in polynomial time."
% =====================================================================================================================
Given an N3DM instance, put $M=3B$ and let $L$ be the least prime with $L\equiv2\pmod3$ and $T:=L-1>12B$. By the prime number theorem
for arithmetic progressions, $L=O(B)$ with an absolute implied constant, which need not be known: $L$ is found by testing the candidates
$12B+2,12B+3,\dots$ by trial division, each test taking time polynomial in $B$. (The variant of Theorem~\ref{thm:conjp} with $L$ a power
of two avoids primes altogether.)

% =====================================================================================================================
% P-M3b  Proof of Theorem thm:conjp (lines 308–310).
% ORIGINAL: "(alternatively the least prime $L$ with $L-1\ge T_0$ and $p\nmid L-1$, $g$ a primitive root, $T=L-1$). In both cases
%   $\ZZ_L^{*}$ has no element of order $p$: for $L=2^{r}$ the group has order $2^{r-1}$, and for the prime choice $p\nmid L-1$."
% =====================================================================================================================
(alternatively the least prime $L$ with $L-1\ge T_0$ and $p\nmid L-1$, $g$ a primitive root, $T=L-1$; such an $L$ is $O(T_0)$ by the prime
number theorem for arithmetic progressions, since the primes $\not\equiv1\pmod p$ have Dirichlet density $(p-2)/(p-1)>0$, it is found by
trial division, and $L\ne p$ because $L>T_0>p$). In both cases $\ZZ_L^{\times}$ has no element of order $p$: for $L=2^{r}$ the group has
order $2^{r-1}$, and for the prime choice $p\nmid L-1$.

% =====================================================================================================================
% P-M4  Proof of Theorem thm:main (line 262).
% ORIGINAL: "$\gamma=1,\beta=0$ gives $S-M-B\in[3-4B,-1]$ and $\gamma=1,\beta=2$ gives $S+M-B\in[3B+3,5B]$, never $\equiv0$;"
% =====================================================================================================================
$\gamma=1,\beta=0$ gives $S-M-B\in[3-4B,-B]$ and $\gamma=1,\beta=2$ gives $S+M-B\in[2B+3,5B]$, never $\equiv0$;

% =====================================================================================================================
% P-M5  Section 4 title, labels, and the preamble of 4.3.
% ORIGINAL (line 208): "\section{Conjugacy by an element of order three is NP-complete}\label{sec:hard}"
% =====================================================================================================================
\section{Conjugacy by an element of prime order}\label{sec:hard}
% ORIGINAL (line 210): "\begin{defi}[\textsc{Order-3 Conjugacy}, $\Dzero$]"
\begin{defi}[\textsc{Order-3 Conjugacy}, $\Dzero$]\label{def:dzero}
% ORIGINAL (line 219): "\subsection{Conjugators between blockwise translations}"
\subsection{Conjugators between blockwise translations}\label{sec:block}
% ORIGINAL (lines 270–278): "For a prime $p$ let $\mathrm{Conj}_p$ be the problem: … $\mathrm{Conj}_3$ is the problem $\Dzero$ of
%   Section~\ref{sec:hard}. This subsection classifies $\mathrm{Conj}_p$ over all primes.
%   \paragraph{The blockwise family and the affine criterion for $x^{p}=1$}
%   Fix $L\ge3$ and $k\ge1$, let $a$ consist of $k$ disjoint $L$-cycles … for $\sigma\in\Sym(k)$ and $t\in\ZZ_L^{k}$."
For a prime $p$ let $\mathrm{Conj}_p$ be the problem: given $a,c\in\Sym(N)$, is there $x\in\Sym(N)$ with $x^{p}=1$ and $xax^{-1}=c$?
Since $x=1$ is a solution iff $a=c$, requiring order exactly $p$ or order dividing $p$ gives polynomially equivalent problems; for $p=3$ this
is the problem of Definition~\ref{def:dzero}. This subsection classifies $\mathrm{Conj}_p$ over all primes.

\paragraph{The affine criterion for $x^{p}=1$}
We use the blockwise family of Section~\ref{sec:block} with an arbitrary modulus $L\ge3$ (not necessarily prime): $a$ acts on each of $k$
blocks $B_i\cong\ZZ_L$ as $z\mapsto z+1$ and $c=c(m)$ as $z\mapsto z+m_i$ with $m_i\in\ZZ_L^{\times}$. Lemma~\ref{lem:block} and its proof
hold verbatim: the conjugators of $a$ to $c$ are the maps sending $B_i$ onto $B_{\sigma(i)}$ by $z\mapsto m_{\sigma(i)}z+t_i$,
$\sigma\in\Sym(k)$, $t\in\ZZ_L^{k}$.
% Globally: replace every "\ZZ_L^{*}" in Section 4.3 by "\ZZ_L^{\times}".

% =====================================================================================================================
% P-M6  Proposition prop:three (lines 69–74): remove the duplicated duality (now only in prop:dual).
% ORIGINAL (line 70): "Every one-variable word with three positive occurrences of $x$ is therefore polynomially equivalent to $\Pthree$, and
%   $\Pthree$ is self-dual: $x\,a\,x^{2}=c$ iff $y\,c\,y^{2}=a$ with $y=x^{-1}$."
% ORIGINAL (line 73, last sentence): "For self-duality, $c=xax^{2}$ gives $x^{-1}cx^{-2}=a$."
% =====================================================================================================================
Every one-variable word with three positive occurrences of $x$ is therefore polynomially equivalent to $\Pthree$ (for the symmetry
$(a,c,x)\mapsto(c,a,x^{-1})$ of $\Pthree$ itself see Proposition~\ref{prop:dual}).
% and delete the sentence "For self-duality, $c=xax^{2}$ gives $x^{-1}cx^{-2}=a$." from the proof.

% =====================================================================================================================
% P-M7  Introduction (line 45).
% ORIGINAL: "Length one is trivial or conjugacy; length two is square-root extraction or conjugacy (Proposition~\ref{prop:short})."
% =====================================================================================================================
Length one is a single equation $x^{\pm1}=c'$; length two is square-root extraction or conjugacy (Proposition~\ref{prop:short}).

% =====================================================================================================================
% P-M8a  Remark rem:notproved (line 397).  PROVISIONAL numbers — to be replaced by the author's recomputation (see audit_report M8a).
% ORIGINAL: "and ``rigid'' instances, all of whose solutions have order three, are rare: in an exhaustive census for $N\le8$ they are about
%   $0.8\%$ of the solvable pairs, mostly with a unique solution."
% =====================================================================================================================
and ``rigid'' instances -- solvable pairs all of whose solutions have order exactly three -- are rare and become rarer: in exhaustive
censuses they are $7.1\%$, $6.2\%$, $3.3\%$ and $1.1\%$ of the solvable pairs for $N=5,6,7,8$, and $78\%$ of them have a unique solution at
$N=8$ (\texttt{audit\_checks.py}, \texttt{rigid\_census2}).

% =====================================================================================================================
% P-M8c  Section 5 (line 480).
% ORIGINAL: "(from $m=70$ with the first-free-point fit to $m=89$--$116$ with the most-constrained fits at a $128$-bit target)"
% =====================================================================================================================
(if the rule requires $\TR(m)\ge2^{128}$: from $m=59$--$74$ with the first-free-point fits to $m=89$--$116$ with the most-constrained fits of
Table~\ref{tab:cost}, the lower value of each pair coming from the $m\log m$ model)
% NOTE: 59 = m log m fit (wreath_summary.json lists 128.17 bits at m=59), 74 = linear fit (1.85m-7.53 >= 128). "70" is not reproducible.

% =====================================================================================================================
% P-M8d  (line 176, already included in P-B2): "for $N\le9$" -> "for $3\le N\le9$".
% P-M8f  Statistics paragraph (line 153).
% ORIGINAL: "and the probabilities of exactly one and exactly two solutions are $0.31$--$0.34$ and $0.16$--$0.19$ (Poisson: $0.368$ and $0.184$);"
% =====================================================================================================================
and, for $N=6,\ldots,9$, the probabilities of exactly one and exactly two solutions are $0.31$--$0.34$ and $0.16$--$0.19$ (Poisson: $0.368$
and $0.184$);

% =====================================================================================================================
% P-M8e  Verification records for thm:conj2 and thm:conjp.  Insert after \end{proof} of thm:conj2 (line 377) and after the remark following
% thm:conjp (line 336).  NUMBERS TO BE MATCHED against conj_p_results.json / verify_campaign.json before use.
% =====================================================================================================================
% after thm:conj2:
The criterion of Theorem~\ref{thm:conj2} was checked against brute force over all involutions of $\Sym(N)$ on $1520$ pairs with
$N\le7$ (\texttt{verify\_campaign.py}) and, independently, on all $873$ pairs $(a,c)$ with $N\le6$, one $a$ per cycle type and every $c$
conjugate to $a$, plus $300$ random pairs at $N=7$ (\texttt{audit\_checks.py}): no discrepancy.
% after the remark following thm:conjp:
Lemma~\ref{lem:affinep} was checked against exhaustive enumeration of all $L^{k}k!$ conjugators on $74$ instances with $p\in\{3,5\}$
($10$ of them positive) and on $320$ further instances with $L\in\{4,5,7,8,9\}$, including composite moduli, and $p\in\{2,3\}$; the
case analysis of Theorem~\ref{thm:conjp} was checked arithmetically on every class-count vector and size vector for $p=3$, $B'\le8$,
$p=5$, $B'\le4$ and $p=7$, $B'=2$. The reduction itself was run against brute-force matching for $q\le3$ ($p=3$) and $q=2$ ($p=5$),
$B'\le6$; the statement that every solution has cycle type $p^{N/p}$ rests on the proof of Lemma~\ref{lem:affinep} and could be tested on
few positive instances only.

% =====================================================================================================================
% P-M9  Open problem 4 (line 487).
% ORIGINAL: "the count restricted to solutions of order three is $\#\mathrm P$-hard by the reduction of Theorem~\ref{thm:main}."
% =====================================================================================================================
the count restricted to solutions of order three is at least as hard as counting the solutions of N3DM: in the reduction of
Theorem~\ref{thm:main} every admissible triple admits exactly $L^{2}$ translation vectors, so the number of order-three conjugators is
$L^{2q}$ times the number of matchings.

% =====================================================================================================================
% P-M10  Lemma lem:inv2, proof (lines 357–360).
% ORIGINAL: "(a) is Lemma~\ref{lem:inv1}. (b) $x^{2}$ commutes with $a$ and $c$ (apply the twisted relations twice), so it is a plain
%   automorphism of the transitive $G$-set $O$ and is the identity iff it fixes one point. (c) Direct verification on each half."
% =====================================================================================================================
(a) A twisted isomorphism $\varphi$ with $\varphi(j_0)=v$ satisfies $\varphi(w(a,c)\,j_0)=w(c,a)\,v$ for every word $w$ (induction on the
length of $w$), and $G$ is transitive on $O$, so $\varphi$ is unique. To decide its existence, propagate $\varphi(au)=c\,\varphi(u)$ and
$\varphi(cu)=a\,\varphi(u)$ from $\varphi(j_0)=v$ along the Schreier graph of $O$; every point of $O$ carries one $a$-edge and one $c$-edge,
so $2|O|$ equations are examined. If no conflict occurs the resulting map satisfies $\varphi a=c\varphi$ and $\varphi c=a\varphi$ on $O$,
its image is a non-empty $\langle a,c\rangle$-invariant subset of the orbit $Gv$, hence all of $Gv$, and it is a bijection iff it is
injective, which is checked in $O(|O|)$ time. (b) $x^{2}a=xcx=ax^{2}$ and likewise $x^{2}c=cx^{2}$, so $x^{2}$ is a plain automorphism of the
transitive $G$-set $O$; a plain automorphism fixing a point is the identity. (c) On $O$, $\psi=\varphi$ gives $\psi a=c\psi$; on $O'$,
$\psi=\varphi^{-1}$ and $\varphi c=a\varphi$ give $\varphi^{-1}a=c\varphi^{-1}$, i.e.\ $\psi a=c\psi$; and $\psi^{2}=1$.

% =====================================================================================================================
% P-M11  Theorem thm:psl, proof (lines 189–193).
% ORIGINAL: "Finite certificate: enumerate $\langle a,c_i\rangle$ by breadth-first search recording a word for each element, define $\varphi$
%   by evaluating the word at $(a,c_2)$, and check that $\varphi$ is a well-defined bijective homomorphism preserving cycle types
%   ($168\times168$ products); count the solutions by brute force over $\Sym(7)$."
% =====================================================================================================================
Proof by finite certificate. Enumerate $G_1=\langle a,c_1\rangle$ by breadth-first search, recording for each $g$ one word $w_g$ with
$g=w_g(a,c_1)$, and put $\varphi(g)=w_g(a,c_2)$. Check $\varphi(gh)=\varphi(g)\varphi(h)$ for all $168^{2}$ pairs: this makes $\varphi$ a
homomorphism, and then every word evaluating to $g$ at $(a,c_1)$ evaluates to $\varphi(g)$ at $(a,c_2)$ (induction on the length), so
$\varphi$ does not depend on the chosen words. Check that $\varphi$ is injective, hence an isomorphism onto $G_2=\langle a,c_2\rangle$
($|G_2|=168$), and that $\ctype\varphi(g)=\ctype g$ for all $g$. Then $w(a,c_2)=\varphi(w(a,c_1))$ has the cycle type of $w(a,c_1)$ for
every word $w$; both groups are transitive; and the permutation characters satisfy $\pi_2\circ\varphi=\pi_1$, since the number of fixed
points is read off the cycle type. The solution counts are obtained by brute force over $\Sym(7)$.

% =====================================================================================================================
% P-M12  Lemma lem:cent, proof (line 103), from "For the criterion," to the end of the proof.
% ORIGINAL: "For the criterion, a cube root $(v,\pi)$ of $(r,\rho)$ in $\ZZ_L\wr\Sym(n)$ satisfies $\pi^{3}=\rho$; … conversely such a
%   grouping yields a cube root. (Verified against exhaustive enumeration of $\Cent(a)$ at $(L,n)=(2,5)$, $60/60$ trials, and $(2,6)$, $40/40$.)"
% =====================================================================================================================
For the criterion, write the elements of $\Cent(a)\cong\ZZ_L\wr\Sym(n)$ as pairs $(v,\pi)$, $\pi$ the permutation of the $n$ cycles of $a$
and $v\in\ZZ_L^{n}$ the rotations, and call the sum of the rotations along a cycle of $\pi$ its rotation sum. A cube root $(v,\pi)$ of
$t=(r,\rho)$ satisfies $\pi^{3}=\rho$, so a $\pi$-cycle of length $\ell$ prime to $3$ is a $\rho$-cycle of the same length and a $\pi$-cycle
of length $3\ell$ cubes to three $\rho$-cycles of length $\ell$. Powers of $t$ show that the rotation sum of $t$ along each of these three
$\rho$-cycles equals the rotation sum of $(v,\pi)$ along the $\pi$-cycle; hence the three sums are equal. Conversely: (i) on a
$\rho$-cycle of length $\ell$ prime to $3$ taken as a single $\pi$-cycle, $\pi=\rho^{k}$ there with $3k\equiv1\pmod\ell$, and the condition
$(v,\pi)^{3}=t$ is the linear system $(I+P+P^{2})v=r$ over $\ZZ_L$, $P$ the $\ell\times\ell$ cyclic permutation matrix of $\pi$; since
$\det(I+P+P^{2})=\prod_{\omega^{\ell}=1}(1+\omega+\omega^{2})=\pm3$ for $3\nmid\ell$ (the map $\omega\mapsto\omega^{3}$ permutes the
$\ell$-th roots of unity), a unit of $\ZZ_L$ because $\gcd(L,3)=1$, the system has a unique solution; (ii) for three $\rho$-cycles of length
$\ell$ with equal rotation sums, let $\pi$ interleave them into one $3\ell$-cycle; now $P$ is a $3\ell$-cycle, the kernel of $I+P+P^{2}$ is
$\{v:\ v_i=w_{i\bmod3},\ w_0+w_1+w_2=0\}$ of order $L^{2}$, so its image has order $L^{3\ell-2}$; every vector in the image has equal partial
sums over the three residue classes of positions modulo $3$ (each equals the total sum of $v$), the three $\rho$-cycles are exactly these
residue classes, and there are exactly $L^{3\ell-2}$ vectors with equal partial sums, so the image is this set and $v$ exists. Hence the
$\rho$-cycles of length divisible by $3$ must be grouped in triples of equal rotation sum, which is possible iff each (length, sum) class
has size divisible by $3$, and every such grouping yields a cube root.
\end{proof}
The criterion was checked against the full enumeration of $\Cent(a)$ for $(L,n)\in\{(2,5),(2,6)\}$ ($60/60$ and $40/40$ trials) and for
$(L,n)\in\{(2,4),(4,3),(5,3),(7,3),(4,4),(5,4),(2,6)\}$, every element of $\Cent(a)$ ($70\,800$ elements, \texttt{audit\_checks.py}).

% =====================================================================================================================
% P-M13  Observation obs:cost (line 450).
% ORIGINAL: "the cost is $|\Alt(m)|^{\Theta(1)}$ with exponent about one fifth of brute force,"
% =====================================================================================================================
the cost is $|\Alt(m)|^{\Theta(1)}$ with exponent between one fifth and two fifths of brute force ($\log_2|\Alt(m)|\approx m\log_2m-1.44m$;
coefficients $0.22$ and $0.37$ in Table~\ref{tab:cost}),

% =====================================================================================================================
% P-M14  Statement of Theorem thm:main (line 216) and thm:conjp (line 302).
% ORIGINAL (216): "Hardness holds already when $a$ consists of $k$ disjoint cycles of one prime length $L\equiv2\pmod3$,"
% =====================================================================================================================
Hardness holds already when $a$ consists of $k$ disjoint cycles of one prime length $L\equiv2\pmod3$ (depending on the instance),
% ORIGINAL (302): "Hardness holds already when $a$ consists of $k$ disjoint cycles of one length $L$,"
Hardness holds already when $a$ consists of $k$ disjoint cycles of one length $L$ (depending on the instance),

% =====================================================================================================================
% STYLE patches
% S2  ORIGINAL (line 407): "\begin{rem}" (the remark after lem:twoeq)   ->
\begin{rem}\label{rem:padding}
% and ORIGINAL (line 484): "(Remark after Lemma~\ref{lem:twoeq})"   ->
(Remark~\ref{rem:padding})
% S3  ORIGINAL (line 197): "has $N!/|\Cent(a)|$ values in general"   ->
has at least $N!/|\Cent(a)|$ values
% S4  ORIGINAL (abstract): "so it is the shortest one-variable shape over permutation groups whose complexity is not classified"   ->
so it is a shortest one-variable shape over permutation groups whose complexity is not classified
% S7  ORIGINAL (line 15): "\newcommand{\Dzero}{\mathsf{Conj}_3}"   ->
\newcommand{\Conj}[1]{\mathsf{Conj}_{#1}}\newcommand{\Dzero}{\Conj{3}}
% and replace every "\mathrm{Conj}_p" by "\Conj{p}", "\mathrm{Conj}_2" by "\Conj{2}", "\mathrm{Conj}_3" by "\Dzero".
% S8  ORIGINAL (line 112): "found no other identical solution under any of the order conditions"   ->
found no other word that is identically a solution under any of the order conditions
% S9  ORIGINAL (line 366): "one (equivalently every) orbit of $\mathcal C$ admits an involutive twisted automorphism."   ->
one (equivalently every) orbit of $\mathcal C$ admits an involutive twisted automorphism, i.e.\ a twisted isomorphism $x\colon O\to O$ with $x^{2}=1$.
% S11 ORIGINAL (line 66): "which exists iff the cycles of even length of $c'c_1$ occur an even number of times for every length"   ->
which exists iff, for every even $m$, the number of $m$-cycles of $c'c_1$ is even
% S12 ORIGINAL (line 201): "the mechanisms of Theorem~\ref{thm:main} are polynomial here"   ->
the mechanisms of Theorem~\ref{thm:main} are polynomial-time here
% S1  Add "(Figure~\ref{fig:red})" after "The reduction" in the first sentence of the proof of thm:main (line 248), "(Figure~\ref{fig:stat})"
%     after "inside a block the solution count behaves like a Poisson variable" (line 153), and "(Figure~\ref{fig:wreath})" after
%     "the single-round system \emph{is} $\Pthree$ inside the wreath product" (line 47 or at thm:wreath, line 436).
